\chapter{Eigenschaften von Markov-Ketten}\label{app:markov-properties}

Markov-Ketten und ihre Zustände lassen sich durch verschiedene Eigenschaften charakterisieren. Diese Eigenschaften helfen, das Langzeitverhalten der Kette besser zu verstehen.

\begin{itemize}
    \item \textbf{Aperiodizität:} Ein Zustand $i$ hat die Periode $k \geq 1$, wenn jede Rückkehr zu $i$ nur in einer Anzahl von Schritten möglich ist, die durch $k$ teilbar ist. Ist $k=1$, so heisst der Zustand \emph{aperiodisch}. Sind alle Zustände aperiodisch, so ist die Markov-Kette aperiodisch.
    \item \textbf{Irreduzibilität:} Eine Markov-Kette heisst \emph{irreduzibel}, wenn es zwischen jedem Zustandspaar $i$ und $j$ eine Kette von Schritten mit positiver Wahrscheinlichkeit gibt.
    \item \textbf{Absorbierende Zustände:} Ein Zustand $i$ ist \emph{absorbierend}, wenn $P_{i,i} = 1$ gilt, d.h. die Kette bleibt nach Erreichen dieses Zustands immer dort.
    \item \textbf{Rekurrenz und Transienz:} Ein Zustand ist \emph{rekurrent}, wenn die Kette mit Wahrscheinlichkeit $1$ irgendwann wieder in diesen Zustand zurückkehrt. Andernfalls ist er \emph{transient}. Ist die erwartete Rückkehrzeit endlich, heisst der Zustand \emph{positiv rekurrent}, sonst \emph{null-rekurrent}.
    \item \textbf{Ergodizität:} Ein Zustand ist \emph{ergodisch}, wenn er positiv rekurrent und aperiodisch ist. Eine Markov-Kette ist ergodisch, wenn alle ihre Zustände ergodisch sind.
\end{itemize}

\begin{mychallenge}[Aperiodizität]
    Gegeben sei die folgende Übergangsmatrix:
    \[
        P = \begin{pmatrix}
            0{,}5 & {,}5 \\
            {,}5  & {,}5 \\
        \end{pmatrix}
    \]
    Untersuchen Sie, ob die Markov-Kette aperiodisch ist.

    \begin{myanswer}
        Von jedem Zustand kann man in einem Schritt in jeden anderen Zustand (auch in sich selbst) gelangen. Die Rückkehr zu einem Zustand ist also in jedem Schritt möglich, insbesondere in aufeinanderfolgenden Schritten. Daher ist die Periode $k=1$ und die Kette ist aperiodisch.
    \end{myanswer}
\end{mychallenge}

\begin{mychallenge}[Irreduzibilität]
    Gegeben sei die Übergangsmatrix:
    \[
        P = \begin{pmatrix}
            0 & 1 \\
            1 & 0 \\
        \end{pmatrix}
    \]
    Ist diese Markov-Kette irreduzibel? Ist sie aperiodisch?

    \begin{myanswer}
        Von jedem Zustand kann man in einem Schritt zum jeweils anderen Zustand gelangen, also ist die Kette irreduzibel. Die Rückkehr zu einem Zustand ist aber nur in einer geraden Anzahl von Schritten möglich (2, 4, 6, ...), daher ist die Periode $k=2$ und die Kette ist periodisch (nicht aperiodisch).
    \end{myanswer}
\end{mychallenge}

\begin{mychallenge}[Absorbierende Zustände]
    Gegeben sei die Übergangsmatrix:
    \[
        P = \begin{pmatrix}
            1    & 0    \\
            {,}3 & {,}7 \\
        \end{pmatrix}
    \]
    Gibt es einen absorbierenden Zustand?

    \begin{myanswer}
        Der Zustand 1 ist absorbierend, da $P_{1,1} = 1$. Das heisst, wenn die Kette einmal in Zustand 1 ist, bleibt sie dort.
    \end{myanswer}
\end{mychallenge}

\begin{mychallenge}[Rekurrenz und Transienz]
    Betrachten Sie die folgende Übergangsmatrix:
    \[
        P = \begin{pmatrix}
            0{,}6 & {,}4 \\
            {,}2  & {,}8 \\
        \end{pmatrix}
    \]
    Ist die Kette irreduzibel? Sind die Zustände rekurrent oder transient?

    \begin{myanswer}
        Von jedem Zustand kann man mit positiver Wahrscheinlichkeit in den anderen gelangen, also ist die Kette irreduzibel. Da es nur zwei Zustände gibt und man immer wieder zurückkehren kann, sind beide Zustände rekurrent.
    \end{myanswer}
\end{mychallenge}

\begin{mychallenge}[Typen von Rekurrenz und Transienz]
    Ordnen Sie für jeden Zustand der folgenden Markov-Ketten die Begriffe \emph{positiv rekurrent}, \emph{null-rekurrent} oder \emph{transient} zu. Begründen Sie Ihre Antwort jeweils kurz.

    \begin{enumerate}
        \item \textbf{Kette A:} \\
              $P = \begin{pmatrix}
                      0{,}5 & {,}5 \\
                      {,}5  & {,}5 \\
                  \end{pmatrix}$

        \item \textbf{Kette B:} \\
              $P = \begin{pmatrix}
                      1    & 0    \\
                      {,}3 & {,}7 \\
                  \end{pmatrix}$

        \item \textbf{Kette C:} (Zustände $1,2,3,\ldots$; von $i$ geht es immer mit Wahrscheinlichkeit $1$ zu $i+1$)

        \item \textbf{Kette D:} (Zustände $1,2,3,\ldots$; von $i$ geht es mit Wahrscheinlichkeit $1/2$ zu $i+1$ und mit $1/2$ zu $i-1$, wobei von $1$ aus nur nach $2$ möglich ist)
    \end{enumerate}

    \begin{myanswer}
        \begin{enumerate}
            \item \textbf{Kette A:} Beide Zustände sind \emph{positiv rekurrent}, da die Rückkehrzeiten endlich sind und die Kette endlich und irreduzibel ist.
            \item \textbf{Kette B:} Zustand 1 ist \emph{positiv rekurrent} (absorbierend), Zustand 2 ist \emph{transient}, da man von 2 mit positiver Wahrscheinlichkeit nach 1 gelangt und dann nie zurückkehrt.
            \item \textbf{Kette C:} Alle Zustände sind \emph{transient}, da man nie in einen Zustand zurückkehren kann (es geht immer nur vorwärts).
            \item \textbf{Kette D:} Alle Zustände sind \emph{null-rekurrent}, da man zwar mit Wahrscheinlichkeit 1 zurückkehrt, aber die erwartete Rückkehrzeit unendlich ist.
        \end{enumerate}
    \end{myanswer}
\end{mychallenge}

\begin{mychallenge}[Ergodizität]
    Welche Bedingungen müssen erfüllt sein, damit eine Markov-Kette ergodisch ist?

    \begin{myanswer}
        Alle Zustände müssen positiv rekurrent und aperiodisch sein.
    \end{myanswer}
\end{mychallenge}

\section{Stationäre Verteilung}

Bei vielen Markov-Ketten interessiert uns besonders das Langzeitverhalten: Gibt es eine stabile Wahrscheinlichkeitsverteilung, die sich nach vielen Schritten einstellt?

\begin{mydefinition}[Stationäre Verteilung]
    Eine Wahrscheinlichkeitsverteilung $\pi = (\pi_1, \pi_2, \ldots)$ heisst \textbf{stationäre Verteilung} einer Markov-Kette mit Übergangsmatrix $P$, wenn gilt:
    \[
        \pi = \pi \cdot P
    \]
    Das bedeutet: Wenn die Kette die Verteilung $\pi$ erreicht hat, bleibt diese Verteilung bei jedem weiteren Schritt unverändert.
\end{mydefinition}

Die stationäre Verteilung kann durch Lösen des linearen Gleichungssystems $\pi = \pi \cdot P$ zusammen mit der Normierungsbedingung $\sum_i \pi_i = 1$ bestimmt werden.

Für endliche, irreduzible und aperiodische Markov-Ketten (siehe \cref{app:markov-properties}) gilt der \textbf{Ergodensatz}: Es existiert genau eine stationäre Verteilung, und die Verteilung nach $n$ Schritten konvergiert gegen diese stationäre Verteilung, unabhängig vom Anfangszustand:
\[
    \lim_{n\to\infty} \pi^{(0)} \cdot P^n = \pi
\]

\begin{myexample}[Stationäre Verteilung beim Wetterbeispiel]
    Betrachten wir nochmals die Übergangsmatrix aus dem Wetterbeispiel:
    \[
        P =
        \begin{pmatrix}
            0.8 & 0.2 \\
            0.4 & 0.6 \\
        \end{pmatrix}
    \]
    Um die stationäre Verteilung $\pi = (\pi_1, \pi_2)$ zu finden, müssen wir das Gleichungssystem $\pi = \pi \cdot P$ lösen:
    \[
        \begin{aligned}
            (\pi_1, \pi_2) & = (\pi_1, \pi_2) \cdot
            \begin{pmatrix}
                0.8 & 0.2 \\
                0.4 & 0.6 \\
            \end{pmatrix}                                            \\
                           & = (0.8\pi_1 + 0.4\pi_2, 0.2\pi_1 + 0.6\pi_2)
        \end{aligned}
    \]

    Dies führt zu den Gleichungen:
    \begin{align}
        \pi_1 & = 0.8\pi_1 + 0.4\pi_2 \\
        \pi_2 & = 0.2\pi_1 + 0.6\pi_2
    \end{align}

    Aus der ersten Gleichung erhalten wir:
    \begin{align}
        \pi_1 - 0.8\pi_1 & = 0.4\pi_2 \\
        0.2\pi_1         & = 0.4\pi_2 \\
        \pi_1            & = 2\pi_2
    \end{align}

    Mit der Normierungsbedingung $\pi_1 + \pi_2 = 1$ ergibt sich:
    \begin{align}
        2\pi_2 + \pi_2 & = 1           \\
        3\pi_2         & = 1           \\
        \pi_2          & = \frac{1}{3}
    \end{align}

    Damit ist $\pi_1 = \frac{2}{3}$. Die stationäre Verteilung ist also $\pi = (\frac{2}{3}, \frac{1}{3})$.

    Das bedeutet: Langfristig ist es an zwei von drei Tagen sonnig und an einem von drei Tagen regnerisch, unabhängig davon, mit welchem Wetter wir starten.
\end{myexample}

\begin{myexercise}
    Bestimmen Sie die stationäre Verteilung für die Markov-Kette mit der Übergangsmatrix
    \[
        P = \begin{pmatrix}
            0.6 & 0.4 \\
            0.1 & 0.9 \\
        \end{pmatrix}
    \]
    aus dem App-Nutzungs-Beispiel.

    \begin{myanswer}
        Wir lösen das Gleichungssystem $\pi = \pi \cdot P$:
        \begin{align}
            (\pi_1, \pi_2) & = (\pi_1, \pi_2) \cdot
            \begin{pmatrix}
                0.6 & 0.4 \\
                0.1 & 0.9 \\
            \end{pmatrix}
        \end{align}

        Dies führt zu:
        \begin{align}
            \pi_1 & = 0.6\pi_1 + 0.1\pi_2 \\
            \pi_2 & = 0.4\pi_1 + 0.9\pi_2
        \end{align}

        Aus der ersten Gleichung:
        \begin{align}
            \pi_1 - 0.6\pi_1 & = 0.1\pi_2 \\
            0.4\pi_1         & = 0.1\pi_2 \\
            4\pi_1           & = \pi_2
        \end{align}

        Mit $\pi_1 + \pi_2 = 1$ ergibt sich:
        \begin{align}
            4\pi_1 + \pi_1 & = 1   \\
            5\pi_1         & = 1   \\
            \pi_1          & = 0.2
        \end{align}

        Also ist $\pi_2 = 0.8$ und die stationäre Verteilung ist $\pi = (0.2, 0.8)$.

        Das bedeutet, dass der Schüler langfristig zu 80\% TikTok und zu 20\% Instagram nutzt, unabhängig davon, mit welcher App er startet.
    \end{myanswer}
\end{myexercise}

\section{Grenzen von Markov-Ketten}

Markov-Ketten sind ein mächtiges Werkzeug zur Modellierung stochastischer Prozesse, aber sie haben auch klare Grenzen:

\begin{itemize}
    \item \textbf{Markov-Eigenschaft:} Die Annahme, dass der nächste Zustand nur vom aktuellen Zustand abhängt, ist in vielen realen Prozessen nicht erfüllt. Oft spielen weiter zurückliegende Zustände oder externe Einflüsse eine Rolle.
    \item \textbf{Feste Übergangswahrscheinlichkeiten:} In klassischen Markov-Ketten sind die Übergangswahrscheinlichkeiten konstant. In der Realität können sie sich jedoch im Zeitverlauf ändern (z.\,B. saisonale Effekte, Lernprozesse).
    \item \textbf{Begrenzte Modellierung komplexer Abhängigkeiten:} Markov-Ketten können keine komplexen, langfristigen Abhängigkeiten oder Rückkopplungen abbilden.
    \item \textbf{Zustandsraum:} Bei sehr vielen oder kontinuierlichen Zuständen wird die Modellierung und Berechnung schnell unübersichtlich oder unmöglich.
    \item \textbf{Keine Kausalität:} Markov-Ketten beschreiben nur Wahrscheinlichkeiten für Übergänge, aber keine Ursachen für Zustandsänderungen.
\end{itemize}

In solchen Fällen sind erweiterte Modelle wie \emph{höherer Ordnung} Markov-Ketten, \emph{Hidden Markov Models} oder andere stochastische Prozesse notwendig.
